\subsubsection{波动方程}\label{subsec:wave-equation}

下一个初等方程是\emph{波动方程}. 这里要寻找依赖于空间变量 $\bm{x}\in\mathbb{R}^n$ 和时间变量 $t\in\mathbb{R}$ 的实值函数 $u$, 使其满足
\begin{equation}\label{eq:wave-general}
u_{tt}=\Delta u.
\end{equation}
同样, Laplace 算子只作用于空间变量. 这一方程描述多种弹性波和电磁波.

再次描述时空柱体 $\Omega_t^+:=\{(\bm{x},t)\in\Omega\times(0,\infty)\}$ 上的一些典型边界条件, 其中 $\Omega$ 是 $\mathbb{R}^n$ 中的有界区域. 由于波动方程关于时间是二阶的, 通常要指定两个初始条件
\begin{equation}\label{eq:wave-general-initial-position}
u(\bm{x},0)=f(\bm{x}),
\end{equation}

\begin{equation}\label{eq:wave-general-initial-velocity}
u_t(\bm{x},0)=g(\bm{x}).
\end{equation}
在弹性力学问题中, 这相当于指定零时刻的位置和速度. 通常在边界的不同部分指定 Dirichlet 或 Neumann 条件. 在弹性力学应用中, 它们通常分别解释为位移条件和牵引条件.

\paragraph*{用分离变量法求解初边值问题}

首先考虑的初边值问题描述一根两端固定, 给定初始位置和速度的单位长度弦. 问题表述如下. 令 $D^+$ 为上一小节定义的 $(x,t)$ 区域. 求函数 $u:\overline{D^+}\to\mathbb{R}$, 使其在 $(x,t)\in D^+$ 时满足一维 (空间) 波动方程
\begin{equation}\label{eq:wave-one-dimensional}
u_{tt}=u_{xx}
\end{equation}
在 $x\in(0,1)$ 时满足初始条件
\begin{equation}\label{eq:wave-interval-initial-position}
u(x,0)=f(x),
\end{equation}
\begin{equation}\label{eq:wave-interval-initial-velocity}
u_t(x,0)=g(x)
\end{equation}
并在 $t>0$ 时满足 Dirichlet 边界条件
\begin{equation}\label{eq:wave-left-boundary}
u(0,t)=0,
\end{equation}
\begin{equation}\label{eq:wave-right-boundary}
u(1,t)=0
\end{equation}


如果同前面一样使用分离变量法, 就得到下列同时满足波动方程 \eqref{eq:wave-one-dimensional} 和边界条件的解族.
\begin{equation}\label{eq:wave-separated-modes}
u_n(x,t)=(\alpha_n\cos n\pi t+\beta_n\sin n\pi t)\sin n\pi x.
\end{equation}


如果初始条件具有如下形式的 Fourier 展开
\begin{equation}\label{eq:wave-position-fourier-series}
f(x)=\sum_{n=1}^{\infty}A_n\sin n\pi x,
\end{equation}

\begin{equation}\label{eq:wave-velocity-fourier-series}
g(x)=\sum_{n=1}^{\infty}B_n\sin n\pi x,
\end{equation}
则该初边值问题的形式级数解为
\begin{equation}\label{eq:wave-interval-series-solution}
u(x,t):=\sum_{n=1}^{\infty}\left(A_n\cos n\pi t+\dfrac{B_n}{n\pi}\sin n\pi t\right)\sin n\pi x.
\end{equation}


\paragraph*{Cauchy 问题的 D'Alembert 解}

\PDEIndexMark{idx:dalembert-heading}\PDEIndexMark{idx:wave-cauchy-problem}本节考虑一维波动方程的 \emph{Cauchy 问题}. 具体地说, 希望求出实值函数 $u$, 使其在半平面 $(x,t)\in(-\infty,\infty)\times(0,\infty)$ 内满足波动方程 \eqref{eq:wave-one-dimensional}, 并在 $x\in(-\infty,\infty)$ 时满足初始条件
\begin{equation}\label{eq:wave-cauchy-initial-position}
u(x,0)=f(x),
\end{equation}

\begin{equation}\label{eq:wave-cauchy-initial-velocity}
u_t(x,0)=g(x)
\end{equation}


为推导这个问题的解, 首先考察波动方程的两个特殊\emph{行波}解. 设 $F$ 和 $G$ 是 $C^2(\mathbb{R})$ 中的实值函数. 可以看到,
\begin{equation}\label{eq:left-traveling-wave}
u_1(x,t):=F(x+t)
\end{equation}
和
\begin{equation}\label{eq:right-traveling-wave}
u_2(x,t):=G(x-t)
\end{equation}
都满足波动方程. 注意, $u_1$ 只是函数 $F$ 以单位速度向左平移, 而 $u_2$ 则是 $G$ 向右平移. 事实上, 可以证明, 波动方程的任意解都具有形式
\begin{equation}\label{eq:wave-traveling-sum}
u(x,t)=F(x+t)+G(x-t).
\end{equation}


为此, 只需作变量变换
\begin{equation}\label{eq:wave-characteristic-xi}
\xi=x+t,
\end{equation}

\begin{equation}\label{eq:wave-characteristic-tau}
\tau=x-t,
\end{equation}
使得
\begin{equation}\label{eq:wave-transformed-function}
\bar u(\xi,\tau):=u\left(\dfrac{\xi+\tau}{2},\dfrac{\xi-\tau}{2}\right).
\end{equation}


利用链式法则可知, 如果 $u$ 满足波动方程, 那么 $\bar u$ 满足
\begin{equation}\label{eq:wave-mixed-derivative-zero}
\bar u_{\xi\tau}=0.
\end{equation}
由此得
\begin{equation}\label{eq:wave-characteristic-split}
\bar u(\xi,\tau)=F(\xi)+G(\tau).
\end{equation}
变回原来的自变量 $(x,t)$, 就得到 \eqref{eq:wave-traveling-sum}.

现在将这一通解形式用于 Cauchy 问题, 代入初始条件 \eqref{eq:wave-cauchy-initial-position} 和 \eqref{eq:wave-cauchy-initial-velocity}, 得到未知函数 $F$ 和 $G$ 满足的如下方程:
\begin{equation}\label{eq:dalembert-initial-position-relation}
f(x)=F(x)+G(x),
\end{equation}

\begin{equation}\label{eq:dalembert-initial-velocity-relation}
g(x)=F'(x)-G'(x).
\end{equation}


由此得到
\begin{equation}\label{eq:dalembert-f-derivative}
F'(x)=\dfrac12(f'(x)+g(x)),
\end{equation}

\begin{equation}\label{eq:dalembert-g-derivative}
G'(x)=\dfrac12(f'(x)-g(x)).
\end{equation}


对这些方程积分, 再将结果代入 \eqref{eq:wave-traveling-sum}, 就得到 \emph{Cauchy 问题的 D'Alembert 解}
\begin{equation}\label{eq:dalembert-cauchy-solution}
u(x,t):=\dfrac12[f(x+t)+f(x-t)]+\dfrac12\int_{x-t}^{x+t}g(s)\,\mathrm{d}s.
\end{equation}


D'Alembert 解 (更具体地说, \eqref{eq:wave-characteristic-split} 所蕴含的解的形式) 最引人注目的特点之一是, 即使 $f$ 和 $g$ 不连续, 解的公式仍然完全有意义. 这样的"解"将包含 $u$ 中的一个以单位速度向左或向右移动的"跳跃". 这类解的存在不应违背我们对波动方程的直觉, 因为我们会称为不连续的物理波动现象 (例如近岸破碎波和爆炸产生的激波) 每天都在发生. 但解在数学上的性质如何? 在解不可微的点, 怎么能说它满足微分方程呢? 后面各章将更充分地研究这个问题, 尤其将在广义波动方程的背景下得到一些相当详细的答案.

\paragraph*{能量守恒}

本节推导波动方程解的一个称为能量守恒的结果. 这里证明适用于上面定义的两端固定的一维波动方程的版本, 将推广留到后面各章.

\begin{Lemma}
\label{lem:wave-energy-conservation}
设 $u:\overline{D^+}\to\mathbb{R}$ 是波动方程 \eqref{eq:wave-one-dimensional} 的一个 $C^2$ 解, 满足边界条件 \eqref{eq:wave-left-boundary} 和 \eqref{eq:wave-right-boundary}. 那么对任意 $t_1\geq t_0\geq0$, 解 $u$ 满足
\begin{equation}\label{eq:wave-energy-conservation}
\int_0^1u_t^2(x,t_1)+u_x^2(x,t_1)\,\mathrm{d}x=\int_0^1u_t^2(x,t_0)+u_x^2(x,t_0)\,\mathrm{d}x.
\end{equation}

\end{Lemma}

\begin{Proof}
同证明热方程的能量不等式时一样, 首先推导一个微分恒等式. 设 $u$ 满足波动方程. 那么
\begin{equation}\label{eq:wave-equation-tested-by-velocity}
0=(u_{tt}-u_{xx})u_t
\end{equation}

\begin{equation}\label{eq:wave-local-energy-identity}
{}=\dfrac{\mathrm{d}}{\mathrm{d}t}(u_t^2+u_x^2)/2-\dfrac{\mathrm{d}}{\mathrm{d}x}(u_xu_t).
\end{equation}
现在利用它在矩形 $(x,t)\in[0,1]\times[t_0,t_1]$ 上积分, 其中可以任意交换积分次序, 得到:
\[
\begin{aligned}
&\int_0^1u_t^2(x,t_1)+u_x^2(x,t_1)\,\mathrm{d}x
-\int_0^1u_t^2(x,t_0)+u_x^2(x,t_0)\,\mathrm{d}x\\
&\quad=\int_0^1\int_{t_0}^{t_1}\dfrac{\mathrm{d}}{\mathrm{d}t}(u_t^2+u_x^2)\,\mathrm{d}t\,\mathrm{d}x\\
&\quad=\int_{t_0}^{t_1}\int_0^12\dfrac{\mathrm{d}}{\mathrm{d}x}(u_xu_t)\,\mathrm{d}x\,\mathrm{d}t\\
&\quad=\int_{t_0}^{t_1}2u_x(1,t)u_t(1,t)\,\mathrm{d}t-\int_{t_0}^{t_1}2u_x(0,t)u_t(0,t)\,\mathrm{d}t.
\end{aligned}
\]
而边界条件 \eqref{eq:wave-left-boundary} 和 \eqref{eq:wave-right-boundary} 蕴含
\begin{equation}\label{eq:wave-boundary-time-derivatives}
u_t(0,t)=u_t(1,t)=0,
\end{equation}
因此得到 \eqref{eq:wave-energy-conservation}.
\end{Proof}

注意, 我们称为波动方程解的能量的量, 与称为热方程解的能量的量, 在数学上看起来很不相同. 然而, 研究这两个量所用的数学技巧是共同的 (将微分方程乘以解或其导数, 再 (实质上) 分部积分以获得估计). 这种获得偏微分方程解的估计的技巧极为有用, 后面各章将加以推广.

\paragraph*{Problems}

\begin{enumerate}[label=\thesection.\arabic*.,ref=\thesection.\arabic*,start=25,series=book-problems]

\item\label{ex:wave-two-boundary-problems} 用分离变量法求解满足下列边界条件的一维波动方程:\begin{enumerate}[label=(\alph*)]
\item $\begin{aligned}u(x,0)&=x^2-x,\\u_t(x,0)&=0,\\u_x(0,t)&=0,\\u_x(1,t)&=0.\end{aligned}$
\item $\begin{aligned}u(x,0)&=0,\\u_t(x,0)&=\sin\pi x,\\u(0,t)&=0,\\u_x(1,t)&=0.\end{aligned}$
\end{enumerate}

\item\label{ex:wave-cauchy-wellposedness} 对区域 $(x,t)\in(-\infty,\infty)\times(0,\infty)$ 上的 Cauchy 问题 \eqref{eq:wave-one-dimensional}, \eqref{eq:wave-cauchy-initial-position}, \eqref{eq:wave-cauchy-initial-velocity}, 给出适定性的具体定义 (特别要明确问题在什么意义下连续依赖于数据). 推导使问题适定的初始数据条件. 如果所考察的区域改为某个 $0<T<\infty$ 对应的 $(x,t)\in(-\infty,\infty)\times(0,T)$, 结果会有何不同. Hint: 如果对 $x\in(-\infty,\infty)$ 有 $u(x,0)=0$ 且 $u_t(x,0)=\epsilon>0$, 那么 $u$ 随时间增长会变得任意大. 找出保证 $u$ 保持有界的初始数据条件.

\item\label{ex:wave-finite-propagation-region} 设 $f$ 和 $g$ 在区间 $[-1,1]$ 外恒为零. 在 $(-\infty,\infty)\times[0,\infty)$ 的什么区域内, 可以保证 Cauchy 问题的解 $u$ 恒为零.

\item\label{ex:heat-propagation-comparison} 对于热方程, 是否有与上一题类似的结果? Hint: 使用
\[
\chi_{[-1,1]}(x):=\begin{cases}1,&x\in[-1,1]\\0,&x\notin[-1,1]\end{cases}
\]
作为初始数据. 利用 Problem~\ref{ex:heat-step-cauchy-problem} 得到一个解.

\item\label{ex:discontinuous-weak-wave-solutions} 将一维波动方程的弱解定义为函数 $u(x,t)$, 它对每个 $\phi\in C_0^2(\mathbb{R}^2)$ 都满足
\begin{equation}\label{eq:wave-distributional-form}
\int_{-\infty}^{\infty}\int_{-\infty}^{\infty}u(x,t)(\phi_{tt}(x,t)-\phi_{xx}(x,t))\,\mathrm{d}x\,\mathrm{d}t=0
\end{equation}
这里, $C_0^2(\mathbb{R}^2)$ 是 $C^2(\mathbb{R}^2)$ 中所有具有紧支集的函数构成的集合; 即在某个有界集外恒为零的函数.
\begin{enumerate}[label=(\alph*)]
\item 证明, 波动方程的任意强解 (经典 $C^2$ 解) 也都是弱解.
\item 证明, 形如
\begin{equation}\label{eq:heaviside-right-moving-solution}
u(x,t):=H(x-t)
\end{equation}
和
\begin{equation}\label{eq:heaviside-left-moving-solution}
u(x,t):=H(x+t)
\end{equation}
的不连续函数都是波动方程的弱解. 这里 $H$ 是 Heaviside 函数:
\begin{equation}\label{eq:heaviside-function}
H(x):=\begin{cases}0,&x<0\\1,&x\geq0.\end{cases}
\end{equation}
\end{enumerate}

\end{enumerate}
